\documentclass[12pt]{article}
\usepackage[utf8]{inputenc}
\usepackage[T1]{fontenc}
\usepackage{hyperref}
\title{Alutor Language Resources for K-12 Curriculum Development}
\author{Compiled by Claude Sonnet 5.5 (claude-sonnet-5-5) for Andr\'as Kornai}
\date{}
\begin{document}
\maketitle

\section{Introduction}
Alutor (ISO 639-3 alr) is a Chukotko-Kamchatkan language of northern Kamchatka, Russia, with very few remaining speakers. Search found no dedicated Alutor grammar, dictionary or textbook; the resources below are background and scholarly discussion only.

\section{Primary Educational Resources}
Search results state that Alutor is taught in the Vyvenka village school, but no specific teaching materials were found. A Russian-language article on native-language teaching in Kamchatka Krai is at \url{https://journals.rudn.ru/polylinguality/article/view/38830}.

\section{Linguistic and Reference Materials}
\begin{itemize}
\item Wikipedia, Alyutor language: \url{https://wikipedia.com/wiki/ISO_639:alr}
\item Nagayama, paper on Alutor language decline, Acta Slavica Iaponica 25: \url{https://src-h.slav.hokudai.ac.jp/publictn/acta/25/nagayama.pdf}
\item Journal article, \emph{Languages and Folklore of the Indigenous Peoples of Siberia}: \url{https://philology.nsc.ru/journals/ykns_en/article.php?id=700}
\end{itemize}

\section{Religious and Community Resources}
None found.

\section{Digital and Audio Resources}
Omniglot has an introduction to Alyutor writing and the language: \url{https://www.omniglot.com/writing/alyutor.htm}

\section{Curriculum Development Context}
Very little published material exists in accessible form. Fortescue's \emph{Comparative Chukotko-Kamchatkan Dictionary} (2005) covers Alutor comparatively, but no URL for it was seen. Russian-language catalogues and the Vyvenka school would be the next places to look.
\end{document}